\section{The integer optimum in the independent-team model}\label{sec:iid}

This section treats the independent-team model, in which the $q^3$ teams have independent
deficits, each with the law of a sum of three independent $U[0,1]$ variables. It is the
market of Section~\ref{sec:model} with the dependence between teams that share a pair
removed; Section~\ref{sec:transfer} carries the conclusion back to the market.

\begin{theorem}\label{thm:iid}
In the independent-team model, with probability tending to $1$ as $q\to\infty$, every
perfect assignment has deficit at least $\betaint\,q^{1/3}$, where $\betaint=2.302$.
\end{theorem}

The proof is a first-moment computation for the number of perfect assignments of small
deficit, carried out conditionally on the stars of the agents (the numbers of cheap teams
at each agent, recorded in bins of cost). The steps are: an exact formula for the conditional first moment in a
configuration model (Subsection~\ref{subsec:iid-firstmoment}); exponential bounds and the
passage to a limit functional on typical stars (Subsection~\ref{subsec:iid-exponent}); a
concavity property of that functional, which bounds its supremum by its value on a tangent
(Subsection~\ref{subsec:iid-concave}); and the evaluation of one explicit number in
interval arithmetic (Subsection~\ref{subsec:iid-proof}). The analytic part is stated as a
criterion, Proposition~\ref{prop:iid-criterion}, valid for every choice of the parameters;
the computer-assisted part is Proposition~\ref{prop:iid-certificate}.

\subsection{The model, the parameters and the criterion}\label{subsec:iid-model}

\subsubsection*{The model}
In this section the three sectors are written $V_1,V_2,V_3$; each has $q$ agents. A team is
a triple $e\in V_1\times V_2\times V_3$. The $q^3$ team deficits $\Defi(e)$ are
independent, and each is distributed as $U+U'+U''$ with $U,U',U''$ independent and uniform
on $[0,1]$. Two facts about the density $f$ of $\Defi(e)$ are used, and nothing else about
the one-team law:
\begin{align}
f(t)&=\tfrac12 t^2\quad\text{for }0\le t\le 1, \tag{F1}\label{eq:iid-F1}\\
f(t)&\le\tfrac12 t^2\quad\text{for all }t\ge 0. \tag{F2}\label{eq:iid-F2}
\end{align}
Indeed $2f(t)=t^2-3(t-1)^2\le t^2$ on $[1,2]$, $2f(t)=(3-t)^2\le1\le t^2$ on $[2,3]$, and
$f=0$ on $(3,\infty)$.

As in Section~\ref{sec:model}, the rescaled cost of a team is $\cost_e=q^{2/3}\Defi(e)$.
For a perfect assignment $N$ (a set of $q$ pairwise disjoint teams) put
$\cost(N)=\sum_{e\in N}\cost_e$. The deficit of $N$ is $q^{-2/3}\cost(N)$, so $N$ has
deficit below $\beta q^{1/3}$ if and only if $\cost(N)<\beta q$. For $\beta\in\R$ let
\[
Z_\beta=\#\{N:\ N\text{ a perfect assignment with }\cost(N)<\beta q\}.
\]
Recall $\Lam(\mathrm{d}x)=\tfrac12x^2\,\mathrm{d}x$. By \eqref{eq:iid-F1}, for every interval
$I\subset[0,q^{2/3}]$,
\begin{equation}\label{eq:iid-binprob}
\P(\cost_e\in I)=\Lam(I)/q^2 .
\end{equation}
Each agent lies in $q^2$ teams, so the number of teams at a given agent with cost in $I$ is
Binomial$(q^2,\Lam(I)/q^2)$, which converges to Poisson$(\Lam(I))$. This is the origin of
the reference star below.

\subsubsection*{Parameters}
Fix a truncation level $K>0$, a number of bins $B\ge1$, and a tilt $\theta>0$. Put
$\eta=K/B$ and $I_b=((b-1)\eta,b\eta]$ for $1\le b\le B$. A team $e$ is \emph{low of bin
$b$} if $\cost_e\in I_b$, and \emph{high} if $\cost_e>K$ (the event $\cost_e=0$ has
probability zero and is ignored). Put
\[
\lambda_b=\Lam(I_b)=\frac{\eta^3(3b^2-3b+1)}{6},\qquad
\sum_{b=1}^B\lambda_b=\frac{K^3}{6},\qquad
\lambda_{\max}=\max_b\lambda_b=\lambda_B,
\]
\[
\kappa_b=\int_{I_b}\mathrm{e}^{-\theta x}\,\tfrac12x^2\dd x,\qquad
J=\int_K^\infty \mathrm{e}^{-\theta x}\,\tfrac12x^2\dd x
 =\mathrm{e}^{-\theta K}\Bigl(\frac{K^2}{2\theta}+\frac{K}{\theta^2}+\frac1{\theta^3}\Bigr).
\]
Vectors in $\R^{B+1}$ have coordinates indexed by $1,\dots,B,\infty$; for $x,y\in\R^{B+1}$
write $\langle x,y\rangle=\sum_{b}x_by_b+x_\infty y_\infty$. The symbol $\ind$ also denotes
the vector $(1,\dots,1)\in\R^{B+1}$. A \emph{weight vector} is any
$u=(u_1,\dots,u_B,u_\infty)\in\R^{B+1}$; put $w_b=\mathrm{e}^{u_b}$ and
$w_\infty=\mathrm{e}^{u_\infty}$. The \emph{reference star} is a random vector
$D=(D_1,\dots,D_B)$ with independent coordinates, $D_b\sim$ Poisson$(\lambda_b)$; its law
on $\N^B$ is denoted $\gamma$. Put
\[
W=W(u)=w_\infty+\sum_{b=1}^BD_bw_b .
\]
Since $w_\infty\le W\le w_\infty+|D|\max_bw_b$, where $|d|=\sum_bd_b$, the variable $\ln W$
has finite moments of all orders.

\begin{definition}\label{def:iid-psi0}
For $u\in\R^{B+1}$ put
\[
\Phi(u)=\E\ln W,\qquad
\rho_b(u)=\E\Bigl[\frac{w_b}{W+w_b}\Bigr]\ (1\le b\le B),\qquad
\eps_0(u)=\E\Bigl[\frac{w_\infty}{W}\Bigr],
\]
\[
C_b(u)=\mathrm{e}^{-3u_b}\,\frac{\kappa_b}{\lambda_b},\qquad
C_\infty(u)=\mathrm{e}^{-3u_\infty}J,
\]
and
\begin{equation}\label{eq:iid-psi0}
\Psi_0(K,B,\theta,u)=3\,\Phi(u)+\sum_{b=1}^B\lambda_b\ln\bigl((1-\rho_b)^3+C_b\rho_b^3\bigr)
 -3\,\eps_0+C_\infty\eps_0^3 ,
\end{equation}
with $\rho_b=\rho_b(u)$, $\eps_0=\eps_0(u)$, $C_b=C_b(u)$, $C_\infty=C_\infty(u)$.
\end{definition}

\begin{proposition}[The criterion]\label{prop:iid-criterion}
Let $K>0$, $B\ge1$, $\theta>0$, $u\in\R^{B+1}$ and $\beta\in\R$ satisfy
\begin{equation}\label{eq:iid-criterion}
\theta\beta+\Psi_0(K,B,\theta,u)<0 .
\end{equation}
Then, in the independent-team model, $\P(Z_\beta\ge1)\to0$ as $q\to\infty$: with
probability tending to $1$ every perfect assignment has deficit at least $\beta q^{1/3}$.
More precisely, there are $\delta>0$ and $c>0$, depending only on $(K,B,\theta,u,\beta)$,
such that for all large $q$
\[
\P(Z_\beta\ge1)\le\P\bigl(D(G)\notin\mathcal{T}_\delta\bigr)+\mathrm{e}^{-cq},
\qquad \P\bigl(D(G)\notin\mathcal{T}_\delta\bigr)\to0,
\]
where $D(G)$ is the array of binned stars defined in
Subsection~\ref{subsec:iid-firstmoment} and $\mathcal{T}_\delta$ is the set of typical
arrays defined in \eqref{eq:iid-typical}. No rate is tracked for the first term.
\end{proposition}

The proof occupies Subsections~\ref{subsec:iid-firstmoment}--\ref{subsec:iid-concave} and
is concluded in Subsection~\ref{subsec:iid-proof}. The following statement is the only
computer-assisted step of this section.

\begin{proposition}[What the computation certifies]\label{prop:iid-certificate}
Let $K=5$, $B=50$ (so $\eta=1/10$), $\theta=1.6712$, and let $u^\ast\in\R^{51}$ be the
vector whose coordinates $u^\ast_1,\dots,u^\ast_{50}$ are, in this order,
\begin{quote}\small\raggedright
$0.71406$, $0.58839$, $0.44812$, $0.30536$, $0.16226$, $0.01936$, $-0.12321$,
$-0.26544$, $-0.40739$, $-0.54916$, $-0.69086$, $-0.83260$, $-0.97452$, $-1.11672$,
$-1.25934$, $-1.40250$, $-1.54629$, $-1.69084$, $-1.83624$, $-1.98257$, $-2.12991$,
$-2.27832$, $-2.42785$, $-2.57854$, $-2.73041$, $-2.88345$, $-3.03766$, $-3.19302$,
$-3.34949$, $-3.50701$, $-3.66554$, $-3.82501$, $-3.98535$, $-4.14649$, $-4.30837$,
$-4.47092$, $-4.63405$, $-4.79770$, $-4.96183$, $-5.12636$, $-5.29125$, $-5.45647$,
$-5.62197$, $-5.78770$, $-5.95361$, $-6.11975$, $-6.28597$, $-6.45224$, $-6.61880$,
$-6.78546$
\end{quote}
and $u^\ast_\infty=-4.61708$. Then
\begin{equation}\label{eq:iid-psi0-certified}
\Psi_0(5,50,1.6712,u^\ast)\le-3.84751328322 .
\end{equation}
Consequently $\Psi_0\le-3.84751328$, $-\Psi_0/\theta\ge2.30224586$, and
$\theta\beta+\Psi_0\le-0.0004108832$ for $\beta=2.302$.
\end{proposition}

\begin{proof}
Inequality \eqref{eq:iid-psi0-certified} is the output of the interval-arithmetic script
\begin{center}
\texttt{lemmaC/round3/final/certI/certify\_rigorous.py}
\end{center}
run on the weights file \texttt{lemmaC/round3/firstmoment/w\_K5\_e10.json} (paths
relative to the lane folder) with mesh parameter $0.005$. The script
encloses $\Phi(u^\ast)$, the $\rho_b(u^\ast)$ and $\eps_0(u^\ast)$ by quadrature of convex
integrands with rigorous error bounds, evaluates \eqref{eq:iid-psi0} in outward-rounded
arithmetic, and returns the upper endpoint of its enclosure of $\Psi_0$ as an exact binary
rational,
\[
-\,\frac{1162839537350571633950917}{2^{78}}\ \le\ -3.84751328322 .
\]
How the enclosures are obtained, and what is trusted, is described in
Section~\ref{sec:certificates}. The remaining statements follow by exact arithmetic:
\begin{align*}
1.6712\times2.30224586&=3.847513281232\le3.84751328322,\\
1.6712\times2.302-3.84751328322&=-0.00041088322 .\qedhere
\end{align*}
\end{proof}

Note that the bound $-\Psi_0/\theta\ge2.30224586$ uses the eleven-decimal form
\eqref{eq:iid-psi0-certified}; the eight-decimal form $\Psi_0\le-3.84751328$ alone gives
only $-\Psi_0/\theta\ge2.30224585$.

\begin{proof}[Proof of Theorem~\ref{thm:iid}, given
Propositions~\ref{prop:iid-criterion} and~\ref{prop:iid-certificate}]
By Proposition~\ref{prop:iid-certificate}, condition \eqref{eq:iid-criterion} holds for
$K=5$, $B=50$, $\theta=1.6712$, $u=u^\ast$ and $\beta=\betaint=2.302$. By
Proposition~\ref{prop:iid-criterion}, $\P(Z_{\betaint}\ge1)\to0$. On the event
$\{Z_{\betaint}=0\}$ every perfect assignment $N$ has $\cost(N)\ge\betaint q$, that is,
deficit at least $\betaint q^{1/3}$.
\end{proof}

\begin{remark}\label{rem:iid-values}
The same argument gives the conclusion of Theorem~\ref{thm:iid} for every
$\beta<-\Psi_0/\theta$, and so for every $\beta\le2.30224586$. A second run of the same
script (weights file \texttt{lemmaC/\allowbreak round3/\allowbreak firstmoment/\allowbreak
w\_K5\_e25.json}: $K=5$, $B=20$, $\theta=1.6671$) certifies $\Psi_0\le-3.83621828$ and $-\Psi_0/\theta\ge2.30113267$; it is
not used. In floating point, the largest value of $-\Psi_0/\theta$ found over $\theta$ and $u$ at $K=5$
is about $2.30114$ for $B=20$ and $2.30225$ for $B=50$, and the corresponding continuum
variational problem has value about $2.30246$ at $K=5$ and $2.30248$ without truncation.
These floating-point values are not used in any proof; they suggest, without proving it,
that the method does not give a constant above about $2.30248$.
\end{remark}

\begin{remark}\label{rem:iid-laws}
The proof of Proposition~\ref{prop:iid-criterion} uses only \eqref{eq:iid-F1},
\eqref{eq:iid-F2} and the independence of the teams.
\end{remark}

\subsection{The conditional first moment}\label{subsec:iid-firstmoment}

Throughout this subsection and the next, $K$, $B$, $\theta$ and $\beta$ are fixed and
$q^{2/3}\ge K$. Put
\[
p_b=\frac{\lambda_b}{q^2}\quad(1\le b\le B),\qquad p=\sum_{b=1}^Bp_b=\frac{K^3}{6q^2}.
\]
By \eqref{eq:iid-binprob}, $p_b=\P(\cost_e\in I_b)$ and $1-p=\P(\cost_e>K)$; moreover
$p\le1/6$ because $q^2\ge K^3$.

\subsubsection*{Tilting, and the split into low and high teams}
Let $G$ be the \emph{labelled low hypergraph}: the set of pairs $(e,b)$ with
$\cost_e\in I_b$. It is the hypergraph of teams of cost at most $K$, each marked by the bin
of its cost. Write $e\in G$ if $e$ is low, and $b(e)$ for its bin. Conditionally on $G$ the
costs are independent. By \eqref{eq:iid-F1} the density of $\cost_e$ is $x^2/(2q^2)$ on
$[0,q^{2/3}]$, so a low team of bin $b$ has conditional density $x^2/(2\lambda_b)$ on
$I_b$. A high team has conditional density $q^{-2/3}f(q^{-2/3}x)/(1-p)$ on $(K,\infty)$,
which is at most $x^2/(2q^2(1-p))$ by \eqref{eq:iid-F2}. Hence
\begin{equation}\label{eq:iid-laplace}
\E\bigl[\mathrm{e}^{-\theta\cost_e}\,\big|\,G\bigr]=\frac{\kappa_b}{\lambda_b}
\ \text{ if }(e,b)\in G,\qquad
\E\bigl[\mathrm{e}^{-\theta\cost_e}\,\big|\,G\bigr]\le\frac{J}{q^2(1-p)}
\ \text{ if $e$ is high.}
\end{equation}
Since $\ind\{\cost(N)<\beta q\}\le\exp(\theta(\beta q-\cost(N)))$, conditional independence
and \eqref{eq:iid-laplace} give
\begin{equation}\label{eq:iid-tilt}
\E[Z_\beta\mid G]\le\mathrm{e}^{\theta\beta q}\sum_{N}\ \prod_{e\in N\cap G}
\frac{\kappa_{b(e)}}{\lambda_{b(e)}}\cdot\Bigl(\frac{J}{q^2(1-p)}\Bigr)^{|N\setminus G|},
\end{equation}
the sum being over all perfect assignments $N$.

A \emph{matching in $G$} is a set of pairwise disjoint low teams. A \emph{profile} is a
vector $m=(m_1,\dots,m_B)\in\N^B$ with $|m|=\sum_bm_b\le q$; a matching has profile $m$ if
it contains exactly $m_b$ teams of bin $b$ for every $b$. For a profile $m$ put
$r=r(m)=q-|m|$, and let $Y_m(G)$ be the number of matchings in $G$ of profile $m$. There
are at most $(q+1)^B$ profiles.

Group the assignments in \eqref{eq:iid-tilt} according to the matching $N\cap G$. If
$N\cap G$ has profile $m$, then $N\setminus G$ is a perfect assignment of the $r+r+r$
agents not covered by $N\cap G$, and there are $(r!)^2$ perfect assignments of these
agents. Dropping the requirement that the teams of $N\setminus G$ be high (a union bound
over the teams above $K$) gives
\begin{equation}\label{eq:iid-cond}
\E[Z_\beta\mid G]\le\mathrm{e}^{\theta\beta q}\sum_{m}Y_m(G)\,
\prod_{b=1}^B\Bigl(\frac{\kappa_b}{\lambda_b}\Bigr)^{m_b}\,(r!)^2
\Bigl(\frac{J}{q^2(1-p)}\Bigr)^{r},
\end{equation}
the sum being over all profiles. An agent not covered by the matching is said to take the
\emph{high option}. The high option cannot be omitted: an agent whose star is empty below
$K$ is covered by no low team, and every perfect assignment uses a high team at it.

\subsubsection*{Degree arrays and the configuration model}
A \emph{degree array} is an array $D=(D^s_{v,b})$ of nonnegative integers, indexed by
$s\in\{1,2,3\}$, $v\in V_s$ and $1\le b\le B$, such that for every $b$ the column sum
\[
n_b=n_b(D)=\sum_{v\in V_s}D^s_{v,b}
\]
is the same for $s=1,2,3$. Put $n=\sum_bn_b$ and $D^s=(D^s_{v,b})_{v\in V_s,\,b\le B}$.
The degree array $D(G)$ of the labelled low hypergraph has $D^s_{v,b}(G)$ equal to the
number of low teams of bin $b$ that contain the agent $v\in V_s$; its row
$(D^s_{v,b}(G))_{b\le B}$ is the \emph{binned star} of $v$. For a degree array $D$ let
$\mathcal{G}(D)$ be the set of labelled hypergraphs with degree array $D$: sets $g$ of
pairs $(e,b)$, $e$ a team and $1\le b\le B$, in which no team occurs twice (neither twice
with the same bin, which is automatic for a set, nor with two different bins), and with
$D(g)=D$.

\smallskip\noindent(a) The teams are independent, each is low of bin $b$ with probability
$p_b$ and high with probability $1-p$. So every $g\in\mathcal{G}(D)$ has
\begin{equation}\label{eq:iid-uniform}
\P(G=g)=\prod_{b=1}^Bp_b^{n_b}\,(1-p)^{q^3-n},
\end{equation}
which depends on $g$ only through $D$: conditionally on $D(G)=D$, the hypergraph $G$ is
uniform on $\mathcal{G}(D)$.

\smallskip\noindent(b) \emph{Configurations.} For each $s$ and $b$ take half-edges
$(v,b,i)$ with $v\in V_s$ and $1\le i\le D^s_{v,b}$; there are $n_b$ half-edges of bin $b$
in each sector. A \emph{configuration} is a family $c=(\sigma_b,\tau_b)_{b\le B}$ in which
$\sigma_b$ and $\tau_b$ are bijections from the half-edges of bin $b$ in sector~$1$ to the
half-edges of bin $b$ in sectors $2$ and $3$ respectively. There are $\prod_b(n_b!)^2$
configurations. The \emph{triples} of $c$ are the $n$ triples of half-edges
$(x,\sigma_b(x),\tau_b(x))$, $x$ a half-edge of bin $b$ in sector~$1$; such a triple has
bin $b$, and its \emph{team} is the triple of agents carrying its three half-edges. A
configuration is \emph{simple} if the teams of its $n$ triples are pairwise distinct; that
is, no two of its triples have the same team, whether the two triples lie in the same bin
or in two different bins. The \emph{projection} of a simple configuration is the set of
pairs (team, bin) of its triples, an element of $\mathcal{G}(D)$. A configuration that is
not simple projects to no element of $\mathcal{G}(D)$.

Every $g\in\mathcal{G}(D)$ is the projection of exactly $\prod_{s,v,b}D^s_{v,b}!$
configurations. Indeed, a configuration projecting to $g$ is obtained by choosing, for
every $s$, $v\in V_s$ and $b$, a bijection between the $D^s_{v,b}$ teams of $g$ of bin $b$
through $v$ and the half-edges $(v,b,1),\dots,(v,b,D^s_{v,b})$; since the teams of $g$ are
distinct, different choices give different sets of triples, hence different configurations.

\smallskip\noindent(c) For a configuration $c$ and a profile $m$ let $\widetilde Y_m(c)$ be
the number of sets of triples of $c$, containing exactly $m_b$ triples of bin $b$ for every
$b$, whose teams are pairwise disjoint. If $c$ is simple with projection $g$, the triples
of $c$ correspond bijectively to the elements of $g$, and $\widetilde Y_m(c)=Y_m(g)$.

\smallskip
For a profile $m$ and $s\in\{1,2,3\}$, a \emph{choice function} of sector $s$ is a
selection of $m_b$ half-edges of bin $b$ in sector $s$, for every $b$, such that the
$|m|$ selected half-edges are at pairwise distinct agents. Let $N_s(m;D^s)$ be the number
of choice functions. Each agent of $V_s$ either has no selected half-edge or has exactly
one, of some bin $b$, chosen among $D^s_{v,b}$; therefore
\begin{equation}\label{eq:iid-genfn}
N_s(m;D^s)=\bigl[z_1^{m_1}\cdots z_B^{m_B}z_\infty^{\,q-|m|}\bigr]\
\prod_{v\in V_s}\Bigl(z_\infty+\sum_{b=1}^BD^s_{v,b}z_b\Bigr).
\end{equation}

\begin{lemma}[Exact count]\label{lem:iid-count}
For every degree array $D$ and every profile $m$ with $m_b\le n_b$ for all $b$, the sum
over all configurations with degree array $D$ is
\[
\sum_c\widetilde Y_m(c)=\prod_{s=1}^3N_s(m;D^s)\cdot\prod_{b=1}^B\bigl(m_b!\,(n_b-m_b)!\bigr)^2 .
\]
Equivalently, the average of $\widetilde Y_m$ over the $\prod_b(n_b!)^2$ configurations
is
\[
\prod_{s=1}^3N_s(m;D^s)\cdot\prod_{b=1}^B\binom{n_b}{m_b}^{-2}.
\]
\end{lemma}

\begin{proof}
Count pairs $(c,S)$ with $S$ a set of triples of $c$ as in (c). The set $S$ determines, in
each sector, the set of half-edges it uses; because the teams of $S$ are pairwise disjoint,
this is a choice function. Conversely, a triple of choice functions (one per sector)
together with, for every $b$, a bijection from the selected half-edges of bin $b$ in
sector~$1$ to those in sector~$2$ and another to those in sector~$3$, is a set $S$ of
triples of half-edges with pairwise disjoint teams. So there are
$\prod_sN_s(m;D^s)\cdot\prod_b(m_b!)^2$ possible sets $S$. A given $S$ is a set of triples
of exactly those configurations whose bijections extend the ones prescribed by $S$; the
remaining $n_b-m_b$ half-edges of bin $b$ in each sector are matched arbitrarily, in
$\prod_b((n_b-m_b)!)^2$ ways.
\end{proof}

\begin{lemma}[Simplicity conditioning]\label{lem:iid-simple}
For a degree array $D$ put
\[
Q(D)=\prod_{b=1}^Bp_b^{n_b}\,(1-p)^{q^3-n}\cdot
\frac{\prod_{b}(n_b!)^2}{\prod_{s,v,b}D^s_{v,b}!}.
\]
Then for every degree array $D$ and every profile $m$ with $m_b\le n_b$ for all $b$,
\begin{equation}\label{eq:iid-simple}
\E\bigl[Y_m(G);\,D(G)=D\bigr]\le Q(D)\cdot\prod_{s=1}^3N_s(m;D^s)\cdot
\prod_{b=1}^B\binom{n_b}{m_b}^{-2},
\end{equation}
while $Y_m(G)=0$ on $\{D(G)=D\}$ if $m_b>n_b$ for some $b$. Moreover
\begin{equation}\label{eq:iid-Qsum}
\sum_DQ(D)\le\exp\Bigl(\frac{K^6}{36\,q\,(1-p)}\Bigr),
\end{equation}
the sum being over all degree arrays.
\end{lemma}

\begin{proof}
A matching of profile $m$ contains $m_b$ distinct low teams of bin $b$, so $Y_m(G)=0$ if
$m_b>n_b(D(G))$ for some $b$. Let $m_b\le n_b$ for all $b$. By \eqref{eq:iid-uniform}, (b)
and (c),
\begin{align*}
\E\bigl[Y_m(G);\,D(G)=D\bigr]
&=\prod_bp_b^{n_b}(1-p)^{q^3-n}\sum_{g\in\mathcal{G}(D)}Y_m(g)\\
&=\frac{\prod_bp_b^{n_b}(1-p)^{q^3-n}}{\prod_{s,v,b}D^s_{v,b}!}
 \sum_{c\text{ simple}}\widetilde Y_m(c).
\end{align*}
Since $\widetilde Y_m\ge0$, the sum over simple configurations is at most the sum over all
configurations, which is given by Lemma~\ref{lem:iid-count}; this is \eqref{eq:iid-simple}.

For \eqref{eq:iid-Qsum}, fix $(n_b)_{b\le B}\in\N^B$. By the multinomial theorem, for each
$s$ and $b$ the sum of $1/\prod_{v\in V_s}D^s_{v,b}!$ over all $(D^s_{v,b})_{v\in V_s}$
with sum $n_b$ equals $q^{n_b}/n_b!$. Hence the sum of $Q(D)$ over the degree arrays with
column sums $(n_b)$ is
\[
(1-p)^{q^3-n}\prod_bp_b^{n_b}(n_b!)^2\Bigl(\frac{q^{n_b}}{n_b!}\Bigr)^3
=(1-p)^{q^3}\prod_b\frac1{n_b!}\Bigl(\frac{q^3p_b}{1-p}\Bigr)^{n_b}.
\]
Summing over $(n_b)\in\N^B$ gives at most
\[
(1-p)^{q^3}\exp\Bigl(\frac{q^3p}{1-p}\Bigr)\le\exp\Bigl(-q^3p+\frac{q^3p}{1-p}\Bigr)
=\exp\Bigl(\frac{q^3p^2}{1-p}\Bigr),
\]
and $q^3p^2=K^6/(36q)$.
\end{proof}

\begin{remark}\label{rem:iid-simple}
If $\mathcal{G}(D)\neq\emptyset$ then $Q(D)$ is the ratio of $\P(D(G)=D)$ to the
probability that a uniform configuration with degree array $D$ is simple.
Lemma~\ref{lem:iid-simple} says that the conditional first moment given the stars is at
most the configuration-model value divided by the probability of simplicity, and that this
division costs a factor $1+o(1)$ after averaging over $D$. No lower bound on the
probability of simplicity is needed.
\end{remark}

\begin{lemma}[Conditional first moment]\label{lem:iid-first-moment}
Let $q^{2/3}\ge K$. For every set $\mathcal{T}$ of degree arrays,
\begin{multline}\label{eq:iid-first-moment}
\E\bigl[Z_\beta;\,D(G)\in\mathcal{T}\bigr]\le
\exp\Bigl(\frac{K^6}{36\,q\,(1-p)}\Bigr)\,\mathrm{e}^{\theta\beta q}\\
\times\sum_{m}\ \sup_{D\in\mathcal{T}:\ n(D)\ge m}\ \Bigl[\,\prod_{s=1}^3N_s(m;D^s)\cdot
\prod_{b=1}^B\binom{n_b}{m_b}^{-2}\Bigr]\\
\times\prod_{b=1}^B\Bigl(\frac{\kappa_b}{\lambda_b}\Bigr)^{m_b}(r!)^2
\Bigl(\frac{J}{q^2(1-p)}\Bigr)^{r},
\end{multline}
where the sum is over all profiles $m$, $r=q-|m|$, the condition $n(D)\ge m$ means
$n_b(D)\ge m_b$ for all $b$, and the supremum of the empty set is $0$.
\end{lemma}

\begin{proof}
Multiply \eqref{eq:iid-cond} by $\ind\{D(G)=D\}$, take expectations, and sum over
$D\in\mathcal{T}$. For each profile $m$ this gives the factor
$\sum_{D\in\mathcal{T}}\E[Y_m(G);D(G)=D]$, which by \eqref{eq:iid-simple} is at most
$\sum_DQ(D)$ times the supremum, over $D\in\mathcal{T}$ with $n(D)\ge m$, of the bracket.
Apply \eqref{eq:iid-Qsum}.
\end{proof}

\subsection{Exponential bounds and the limit functional}\label{subsec:iid-exponent}

\subsubsection*{Choice functions}
For a degree array $D$, $s\in\{1,2,3\}$ and $u\in\R^{B+1}$ put
\[
\widehat\Phi_s(u)=\widehat\Phi_s(u;D)=\frac1q\sum_{v\in V_s}
\ln\Bigl(\mathrm{e}^{u_\infty}+\sum_{b=1}^BD^s_{v,b}\,\mathrm{e}^{u_b}\Bigr).
\]
For a profile $m$ write $\mu_b=m_b/q$, $\eps=r/q=1-\sum_b\mu_b$, and
$\langle\mu,u\rangle=\sum_b\mu_bu_b+\eps u_\infty$. The polynomial in \eqref{eq:iid-genfn}
has nonnegative coefficients, so each coefficient is at most the value of the polynomial
at a positive point divided by the corresponding monomial. Taking $z_b=\mathrm{e}^{u_b}$,
$z_\infty=\mathrm{e}^{u_\infty}$ gives, for every $u\in\R^{B+1}$ and every profile $m$,
\begin{equation}\label{eq:iid-choice}
N_s(m;D^s)\le\exp\Bigl(q\bigl[\widehat\Phi_s(u)-\langle\mu,u\rangle\bigr]\Bigr).
\end{equation}

\subsubsection*{Binomials and factorials}
Let $h(x)=-x\ln x-(1-x)\ln(1-x)$ be the binary entropy, and $g(x,y)=y\,h(x/y)$ for
$0\le x\le y$, $y>0$, with $g(0,0)=0$; then $0\le g(x,y)\le y\ln2$ and $g$ is continuous on
$\{0\le x\le y\}$. The following elementary bounds are used:
$\binom nm\ge\exp(n\,h(m/n))/(n+1)$ for $0\le m\le n$, and
$r!\le\mathrm{e}\,(r+1)(r/\mathrm{e})^r$ for $r\ge0$ (with $0^0=1$).

\begin{lemma}[The exponent]\label{lem:iid-exponent}
Let $q^{2/3}\ge K$. For a degree array $D$ and a profile $m$ with $m_b\le n_b(D)$ for all
$b$, put $\nu_b=n_b/q$ and
\begin{multline}\label{eq:iid-Xi}
\Xi(\mu;D)=\theta\beta+\sum_{s=1}^3\inf_{u\in\R^{B+1}}
\bigl[\widehat\Phi_s(u;D)-\langle\mu,u\rangle\bigr]-2\sum_{b=1}^Bg(\mu_b,\nu_b)\\
+\sum_{b=1}^B\mu_b\ln\frac{\kappa_b}{\lambda_b}+\eps\ln J+2\eps\ln\eps-2\eps
\end{multline}
(with $0\ln0=0$). Then for every set $\mathcal{T}$ of degree arrays
\begin{equation}\label{eq:iid-exponent}
\E\bigl[Z_\beta;\,D(G)\in\mathcal{T}\bigr]\le c_q\,(q+1)^{B+2}\
\sup\ \Bigl[\,\prod_{b=1}^B(n_b+1)^2\cdot\exp\bigl(q\,\Xi(\mu;D)\bigr)\Bigr],
\end{equation}
where the supremum is over all $D\in\mathcal{T}$ and all profiles $m$ with $m_b\le n_b(D)$
for all $b$, and
\[
c_q=\mathrm{e}^2\exp\Bigl(\frac{K^6}{36\,q\,(1-p)}\Bigr)(1-p)^{-q}
\le\mathrm{e}^2\exp\Bigl(\frac{K^{9/2}}{30}+\frac{K^{3/2}}{5}\Bigr).
\]
The factor $\prod_b(n_b+1)^2$ depends on the degree array and stays inside the supremum.
\end{lemma}

\begin{proof}
Fix $D$ and $m$ with $m_b\le n_b$. By \eqref{eq:iid-choice}, optimised over $u$ separately
for each $s$,
$\prod_sN_s(m;D^s)\le\exp(q\sum_s\inf_u[\widehat\Phi_s(u)-\langle\mu,u\rangle])$. By the
bound on binomial coefficients,
$\prod_b\binom{n_b}{m_b}^{-2}\le\prod_b(n_b+1)^2\exp(-2n_bh(m_b/n_b))
=\prod_b(n_b+1)^2\cdot\exp(-2q\sum_bg(\mu_b,\nu_b))$. Next,
$\prod_b(\kappa_b/\lambda_b)^{m_b}=\exp(q\sum_b\mu_b\ln(\kappa_b/\lambda_b))$. Finally, by
the bound on $r!$ and $r=\eps q$,
\begin{align*}
(r!)^2\Bigl(\frac{J}{q^2(1-p)}\Bigr)^r
&\le\mathrm{e}^2(r+1)^2(1-p)^{-r}\Bigl(\frac{r}{\mathrm{e}q}\Bigr)^{2r}J^r\\
&\le\mathrm{e}^2(q+1)^2(1-p)^{-q}\exp\bigl(q[2\eps\ln\eps-2\eps+\eps\ln J]\bigr).
\end{align*}
So the term of \eqref{eq:iid-first-moment} belonging to $m$ and $D$, times
$\mathrm{e}^{\theta\beta q}$, is at most
$\mathrm{e}^2(q+1)^2(1-p)^{-q}\prod_b(n_b+1)^2\exp(q\,\Xi(\mu;D))$. There are at most
$(q+1)^B$ profiles, which gives \eqref{eq:iid-exponent}. For the bound on $c_q$ use
$p\le1/6$ and $q\ge K^{3/2}$: then $K^6/(36q(1-p))\le K^{9/2}/30$ and
$(1-p)^{-q}\le\exp(qp/(1-p))\le\exp(K^3/(5q))\le\exp(K^{3/2}/5)$.
\end{proof}

\subsubsection*{Typical stars}
For $s\in\{1,2,3\}$ let $\widehat\gamma_s$ be the empirical law on $\N^B$ of the $q$ binned
stars of sector $s$: $\widehat\gamma_s(d)=q^{-1}\#\{v\in V_s:(D^s_{v,b})_{b\le B}=d\}$.
Recall that $\gamma$ is the law of the reference star and $|d|=\sum_bd_b$. For $\delta>0$
let $\mathcal{T}_\delta$ be the set of degree arrays $D$ with
\begin{equation}\label{eq:iid-typical}
\sum_{d\in\N^B}\bigl|\widehat\gamma_s(d)-\gamma(d)\bigr|\,(1+|d|)\le\delta
\qquad\text{for }s=1,2,3 .
\end{equation}

\begin{lemma}[Stars are typical]\label{lem:iid-typical}
For every $\delta>0$, $\P(D(G)\in\mathcal{T}_\delta)\to1$ as $q\to\infty$.
\end{lemma}

\begin{proof}
Fix $s$. The binned star of $v\in V_s$ is a function of the costs of the $q^2$ teams
containing $v$, and these sets of teams are disjoint for distinct agents of $V_s$. So the
$q$ stars of sector $s$ are independent, with common law $\gamma_q$, the law of
$\sum_{i\le q^2}\zeta_i$ where the $\zeta_i$ are independent, $\zeta_i=\mathbf{e}_b$ (the
$b$th unit vector of $\N^B$) with probability $p_b$ and $\zeta_i=0$ with probability
$1-p$.

\emph{Total variation.} Let $\zeta'$ have independent Poisson$(p_b)$ coordinates. Then
$\P(\zeta'=0)=\mathrm{e}^{-p}\ge1-p$ and
$\P(\zeta'=\mathbf{e}_b)=p_b\mathrm{e}^{-p}\le p_b$, so the total variation distance
between the laws of $\zeta_1$ and $\zeta'$ is
$\sum_b(p_b-p_b\mathrm{e}^{-p})=p(1-\mathrm{e}^{-p})\le p^2$. The sum of $q^2$ independent
copies of $\zeta'$ has law $\gamma$, because $q^2p_b=\lambda_b$. Coupling term by term,
\[
d_{\mathrm{TV}}(\gamma_q,\gamma)\le q^2p^2=\frac{K^6}{36q^2}.
\]

\emph{Small stars.} Let $R\ge1$. For each $d$, $q\widehat\gamma_s(d)$ is
Binomial$(q,\gamma_q(d))$, so
$\E|\widehat\gamma_s(d)-\gamma_q(d)|\le(\gamma_q(d)/q)^{1/2}$. Hence
\begin{align*}
\E\sum_{|d|\le R}\bigl|\widehat\gamma_s(d)-\gamma(d)\bigr|(1+|d|)
&\le(1+R)\Bigl[\sum_{|d|\le R}\Bigl(\frac{\gamma_q(d)}{q}\Bigr)^{1/2}
 +2\,d_{\mathrm{TV}}(\gamma_q,\gamma)\Bigr]\\
&\le(1+R)\Bigl[\frac{(R+1)^B}{q^{1/2}}+\frac{K^6}{18q^2}\Bigr],
\end{align*}
because there are at most $(R+1)^B$ vectors $d\in\N^B$ with $|d|\le R$ and
$\sum_d|\gamma_q(d)-\gamma(d)|=2\,d_{\mathrm{TV}}(\gamma_q,\gamma)$. The right side tends to
$0$ as $q\to\infty$ for fixed $R$.

\emph{Large stars.} Under $\gamma$ the total $|D|$ is Poisson$(K^3/6)$, and under
$\gamma_q$ it is Binomial$(q^2,p)$, with the same mean $K^3/6$ and variance at most
$K^3/6$. In both cases $\E(1+|D|)^2\le C:=(1+K^3/6)^2+K^3/6$, and so
$\E[(1+|D|);|D|>R]\le C/(1+R)$. Therefore
\begin{multline*}
\E\sum_{|d|>R}\bigl|\widehat\gamma_s(d)-\gamma(d)\bigr|(1+|d|)\\
\le\E_{\gamma_q}\bigl[(1+|D|);|D|>R\bigr]+\E_{\gamma}\bigl[(1+|D|);|D|>R\bigr]
\le\frac{2C}{1+R},
\end{multline*}
uniformly in $q$.

Given $\delta>0$ and $\omega>0$, choose $R$ with $2C/(1+R)\le\delta\omega/12$ and then
$q_0$ so large that for $q\ge q_0$ the bound for small stars is at most $\delta\omega/12$.
For $q\ge q_0$ the left side of \eqref{eq:iid-typical} has expectation at most
$\delta\omega/6$, so by Markov's inequality it exceeds $\delta$ with probability at most
$\omega/6$, and a union bound over $s=1,2,3$ gives
$\P(D(G)\notin\mathcal{T}_\delta)\le\omega/2$.
\end{proof}

Two consequences of \eqref{eq:iid-typical} are used. Let $D\in\mathcal{T}_\delta$.
\begin{itemize}
\item[(T1)] $|\nu_b-\lambda_b|\le\delta$ for every $b$, because
$\nu_b=\sum_d\widehat\gamma_s(d)\,d_b$ and $\lambda_b=\sum_d\gamma(d)\,d_b$. In particular
$n_b+1\le1+(\lambda_{\max}+\delta)q$.
\item[(T2)] Let $A>0$ and $u\in[-A,0]^{B+1}$. For every $d\in\N^B$,
$\mathrm{e}^{-A}\le\mathrm{e}^{u_\infty}+\sum_bd_b\mathrm{e}^{u_b}\le1+|d|$, so the
logarithm of the middle term is at most $A+|d|\le(A+1)(1+|d|)$ in absolute value. Since $\widehat\Phi_s(u;D)$ and $\Phi(u)$ are the integrals of
$d\mapsto\ln(\mathrm{e}^{u_\infty}+\sum_bd_b\mathrm{e}^{u_b})$ against $\widehat\gamma_s$
and $\gamma$ respectively,
\begin{equation}\label{eq:iid-phihat}
\bigl|\widehat\Phi_s(u;D)-\Phi(u)\bigr|\le(A+1)\,\delta\qquad
\text{for all }u\in[-A,0]^{B+1},\ s=1,2,3 .
\end{equation}
\end{itemize}

\subsubsection*{The limit functional}
Let
\[
\mathcal{P}=\Bigl\{\mu=(\mu_1,\dots,\mu_B,\eps)\in\R^{B+1}:\ 0\le\mu_b\le\lambda_b,\
\eps\ge0,\ \sum_{b=1}^B\mu_b+\eps=1\Bigr\}.
\]
For $\mu\in\R^{B+1}$ with coordinate sum $1$ (the last coordinate being written $\eps$)
and $A>0$ define
\begin{align*}
H(\mu)&=\inf_{u\in\R^{B+1}}\bigl[\Phi(u)-\langle\mu,u\rangle\bigr],&
H_A(\mu)&=\min_{u\in[-A,0]^{B+1}}\bigl[\Phi(u)-\langle\mu,u\rangle\bigr],
\end{align*}
and for $\mu\in\mathcal{P}$
\begin{align*}
S_0(\mu)&=\sum_{b=1}^B\lambda_b\,h\Bigl(\frac{\mu_b}{\lambda_b}\Bigr)-\eps\ln\eps+\eps,\\
L(\mu)&=\theta\beta+3H(\mu)-2S_0(\mu)+\langle\ell,\mu\rangle,\qquad
\ell_b=\ln\frac{\kappa_b}{\lambda_b},\quad\ell_\infty=\ln J .
\end{align*}
Put $L^\ast=\sup_{\mu\in\mathcal{P}}L(\mu)$. Here $H(\mu)\in[-\infty,\infty)$, and $H$ is
concave, being an infimum of affine functions. Since
$\Phi(u+t\ind)=\Phi(u)+t$ and $\langle\mu,\ind\rangle=1$, the bracket is invariant under
$u\mapsto u+t\ind$; one may therefore restrict the infimum to $u$ with $\max u=0$, and
consequently
\begin{equation}\label{eq:iid-HA}
H_A(\mu)\downarrow H(\mu)\quad\text{as }A\uparrow\infty .
\end{equation}

\begin{lemma}[Passage to the limit functional]\label{lem:iid-limit}
For $\delta>0$ let $\mathcal{K}_\delta$ be the set of pairs $(\mu,\nu)$ with
$\mu=(\mu_1,\dots,\mu_B,\eps)$ a probability vector, $\nu\in\R^B$,
$|\nu_b-\lambda_b|\le\delta$ and $0\le\mu_b\le\nu_b$ for all $b$; and for $A>0$ put
\[
L_A(\mu,\nu)=\theta\beta+3H_A(\mu)-2\sum_{b=1}^Bg(\mu_b,\nu_b)+\langle\ell,\mu\rangle
+2\eps\ln\eps-2\eps .
\]
Then for all $A>0$ and $\delta>0$,
\begin{equation}\label{eq:iid-limsup}
\limsup_{q\to\infty}\frac1q\ln\E\bigl[Z_\beta;\,D(G)\in\mathcal{T}_\delta\bigr]
\le\sup_{\mathcal{K}_\delta}L_A+3(A+1)\,\delta,
\end{equation}
and
\begin{equation}\label{eq:iid-inf}
\inf_{A>0}\ \inf_{\delta>0}\ \Bigl[\sup_{\mathcal{K}_\delta}L_A+3(A+1)\,\delta\Bigr]\le L^\ast .
\end{equation}
In particular, if $L^\ast<0$ there are $\delta>0$ and $c>0$ such that
$\E[Z_\beta;D(G)\in\mathcal{T}_\delta]\le\mathrm{e}^{-cq}$ for all large $q$.
\end{lemma}

\begin{proof}
The function $H_A$ is a minimum of affine functions over a compact set; it is finite and
concave on $\R^{B+1}$, hence continuous. The functions $g$ and $\eps\mapsto\eps\ln\eps$ are
continuous. So $L_A$ is continuous on the compact set $\mathcal{K}_\delta$, and its
supremum is finite.

\emph{Proof of \eqref{eq:iid-limsup}.} Let $D\in\mathcal{T}_\delta$ and let $m$ be a
profile with $m_b\le n_b$ for all $b$. Then $(\mu,\nu)\in\mathcal{K}_\delta$ by (T1). By
\eqref{eq:iid-phihat},
\[
\inf_{u\in\R^{B+1}}\bigl[\widehat\Phi_s(u;D)-\langle\mu,u\rangle\bigr]
\le\min_{u\in[-A,0]^{B+1}}\bigl[\widehat\Phi_s(u;D)-\langle\mu,u\rangle\bigr]
\le H_A(\mu)+(A+1)\,\delta
\]
for each $s$, so that $\Xi(\mu;D)\le L_A(\mu,\nu)+3(A+1)\delta$. By (T1),
$\prod_b(n_b+1)^2\le(1+(\lambda_{\max}+\delta)q)^{2B}$, uniformly in
$D\in\mathcal{T}_\delta$. Lemma~\ref{lem:iid-exponent} with
$\mathcal{T}=\mathcal{T}_\delta$ now gives
\begin{multline*}
\E\bigl[Z_\beta;\,D(G)\in\mathcal{T}_\delta\bigr]\\
\le c_q\,(q+1)^{B+2}\bigl(1+(\lambda_{\max}+\delta)q\bigr)^{2B}
\exp\Bigl(q\Bigl[\sup_{\mathcal{K}_\delta}L_A+3(A+1)\,\delta\Bigr]\Bigr),
\end{multline*}
and the factor in front of the exponential is polynomial in $q$.

\emph{Proof of \eqref{eq:iid-inf}.} Let $c>L^\ast$. On $\mathcal{P}$,
$-2\sum_bg(\mu_b,\lambda_b)+2\eps\ln\eps-2\eps=-2S_0(\mu)$, so
$L_A(\mu,\lambda)=\theta\beta+3H_A(\mu)-2S_0(\mu)+\langle\ell,\mu\rangle$, which by
\eqref{eq:iid-HA} decreases to $L(\mu)$ as $A\uparrow\infty$, for every
$\mu\in\mathcal{P}$. The sets $\{\mu\in\mathcal{P}:L_A(\mu,\lambda)\ge c\}$ are compact and
decrease as $A$ increases; a common point $\mu$ would have $L(\mu)\ge c>L^\ast$, which is
impossible. So one of them is empty: there is $A$ with
$\sup_{\mu\in\mathcal{P}}L_A(\mu,\lambda)<c$. Fix this $A$. The compact sets
$\mathcal{K}_\delta$ decrease, as $\delta\downarrow0$, to
$\mathcal{K}_0=\mathcal{P}\times\{\lambda\}$. The compact sets
$\{(\mu,\nu)\in\mathcal{K}_\delta:L_A(\mu,\nu)\ge c\}$ decrease with $\delta$ and have
empty intersection, so one of them is empty: there is $\delta_0>0$ with
$\sup_{\mathcal{K}_\delta}L_A<c$ for $\delta\le\delta_0$. Letting $\delta\downarrow0$, the
left side of \eqref{eq:iid-inf} is at most $c$. As $c>L^\ast$ was arbitrary,
\eqref{eq:iid-inf} follows. (If $L^\ast=-\infty$ the same argument applies to every real
$c$.)

The last assertion follows from \eqref{eq:iid-limsup} and \eqref{eq:iid-inf}: choose $A$
and $\delta$ with $\sup_{\mathcal{K}_\delta}L_A+3(A+1)\delta<0$, and let $c$ be half the
absolute value of this number.
\end{proof}

The argument is by compactness: it gives no explicit $A$ and $\delta$, and no rate.
Lemmas~\ref{lem:iid-first-moment} and~\ref{lem:iid-exponent} are stated for an arbitrary
set $\mathcal{T}$ of degree arrays, so that they can be applied to sets of arrays other
than $\mathcal{T}_\delta$.

\subsection{Concavity and the tangent inequality}\label{subsec:iid-concave}

This subsection concerns only the reference star $D$ (independent Poisson$(\lambda_b)$
coordinates) and is exact for finite $B$. Its outcome is
Proposition~\ref{prop:iid-tangent-bound}: $L^\ast\le\theta\beta+\Psi_0(K,B,\theta,u_0)$ for
every weight vector $u_0$.

\subsubsection*{Notation}
For $u\in\R^{B+1}$, with $w=\mathrm{e}^u$ and $W=w_\infty+\sum_bD_bw_b$ as before, define
the random probability vector $\pi=\pi(u)$ on $\{1,\dots,B,\infty\}$ by
\[
\pi_b=\frac{D_bw_b}{W}\quad(1\le b\le B),\qquad\pi_\infty=\frac{w_\infty}{W}.
\]
The function $u\mapsto\ln W$ is smooth, with gradient $\pi$ and Hessian
$\operatorname{diag}(\pi)-\pi\pi^{\mathsf T}$; all its partial derivatives of positive
order are polynomials in the coordinates of $\pi$, hence bounded, so $\Phi$ is smooth and
may be differentiated under the expectation:
\[
\nabla\Phi(u)=\E\,\pi=:\mu(u),\qquad
\nabla^2\Phi(u)=\Sigma(u):=\E\bigl[\operatorname{diag}(\pi)-\pi\pi^{\mathsf T}\bigr].
\]
Let $\mathcal{H}_0=\{x\in\R^{B+1}:\langle x,\ind\rangle=0\}$ and
$\mathcal{H}_1=\{x\in\R^{B+1}:\langle x,\ind\rangle=1\}$. Then $\mu(u)\in\mathcal{H}_1$,
$\Sigma(u)\ind=0$, and for $x\in\R^{B+1}$
\[
x^{\mathsf T}\Sigma(u)\,x=\E\Bigl[\sum_i\pi_ix_i^2-\Bigl(\sum_i\pi_ix_i\Bigr)^2\Bigr]\ge0,
\]
the expected variance of $x$ under $\pi$. With positive probability all $D_b\ge1$, and then
$\pi$ charges every coordinate; so $x^{\mathsf T}\Sigma(u)x>0$ unless $x$ is a multiple of
$\ind$. Thus $\Sigma(u)$ maps $\mathcal{H}_0$ to itself and its restriction
$\Sigma(u)|_{\mathcal{H}_0}$ is positive definite; $\Phi$ is convex on $\R^{B+1}$ and
strictly convex on $\mathcal{H}_0$. Since $\Phi(u+t\ind)=\Phi(u)+t$, the quantities
$\mu(u)$, $\Sigma(u)$, $\rho_b(u)$ and $\eps_0(u)$ are unchanged when $u$ is replaced by
$u+t\ind$.

\begin{lemma}[Poisson identities]\label{lem:iid-mecke}
For every $u\in\R^{B+1}$: $\mu_b(u)=\lambda_b\rho_b(u)$ with $\rho_b(u)\in(0,1)$ for
$1\le b\le B$; $\mu_\infty(u)=\eps_0(u)\in(0,1]$; and
\begin{equation}\label{eq:iid-pipi}
\E\bigl[\pi\pi^{\mathsf T}\bigr]=\int_0^\infty t\,\E\bigl[\mathrm{e}^{-tW}\bigr]\,
v_tv_t^{\mathsf T}\dd t+\operatorname{diag}\Bigl(\lambda_b\,
\E\Bigl[\Bigl(\frac{w_b}{W+w_b}\Bigr)^2\Bigr]\ (b\le B),\ 0\Bigr),
\end{equation}
where $v_t\in\R^{B+1}$ has coordinates $(v_t)_b=\lambda_bw_b\mathrm{e}^{-tw_b}$ and
$(v_t)_\infty=w_\infty$.
\end{lemma}

\begin{proof}
Let $\mathbf{e}_b$ be the $b$th unit vector of $\N^B$. For independent Poisson variables
$D_b$ with means $\lambda_b$ and every bounded function $\chi$ on $\N^B$, the Mecke formula
gives
\begin{align*}
\E[D_b\,\chi(D)]&=\lambda_b\,\E\,\chi(D+\mathbf{e}_b),\\
\E[D_bD_c\,\chi(D)]&=\lambda_b\lambda_c\,\E\,\chi(D+\mathbf{e}_b+\mathbf{e}_c)\qquad(b\ne c),\\
\E[D_b^2\,\chi(D)]&=\lambda_b^2\,\E\,\chi(D+2\mathbf{e}_b)+\lambda_b\,\E\,\chi(D+\mathbf{e}_b).
\end{align*}
Adding $\mathbf{e}_b$ to $D$ replaces $W$ by $W+w_b$. The first formula with
$\chi=w_b/W$ gives $\mu_b(u)=\E[D_bw_b/W]=\lambda_b\E[w_b/(W+w_b)]$, and
$0<w_b/(W+w_b)<1$. By definition $\mu_\infty(u)=\E[w_\infty/W]=\eps_0(u)$, and
$0<w_\infty/W\le1$.

For \eqref{eq:iid-pipi} apply the formulas to $\chi=w_bw_c/W^2$, $w_b^2/W^2$ and
$w_bw_\infty/W^2$. Entry by entry, for $b\ne c$,
\begin{align*}
\E[\pi_b\pi_c]&=\lambda_b\lambda_cw_bw_c\,\E\bigl[(W+w_b+w_c)^{-2}\bigr],\\
\E[\pi_b^2]&=\lambda_b^2w_b^2\,\E\bigl[(W+2w_b)^{-2}\bigr]
 +\lambda_bw_b^2\,\E\bigl[(W+w_b)^{-2}\bigr],\\
\E[\pi_b\pi_\infty]&=\lambda_bw_bw_\infty\,\E\bigl[(W+w_b)^{-2}\bigr],\qquad
\E[\pi_\infty^2]=w_\infty^2\,\E\bigl[W^{-2}\bigr].
\end{align*}
Now insert $x^{-2}=\int_0^\infty t\,\mathrm{e}^{-tx}\dd t$ for $x>0$ into every term
except the second term of $\E[\pi_b^2]$, and exchange expectation and integral (the
integrands are nonnegative). For instance
\[
\lambda_b\lambda_cw_bw_c\,\E\bigl[(W+w_b+w_c)^{-2}\bigr]
=\int_0^\infty t\,\E\bigl[\mathrm{e}^{-tW}\bigr]\,(v_t)_b(v_t)_c\dd t .
\]
The four families of terms
are the entries of $\int_0^\infty t\,\E[\mathrm{e}^{-tW}]\,v_tv_t^{\mathsf T}\dd t$, and
the second term of $\E[\pi_b^2]$ is the diagonal matrix in \eqref{eq:iid-pipi}. The
integral converges because $W\ge w_\infty>0$.
\end{proof}

\begin{lemma}[Covariance domination]\label{lem:iid-domination}
For every $u\in\R^{B+1}$, in the positive semidefinite order,
\[
\Sigma(u)\preceq\Gamma(u):=\operatorname{diag}\bigl(\lambda_b\rho_b(1-\rho_b)\ (b\le B),\
\eps_0\bigr),\qquad\rho_b=\rho_b(u),\ \eps_0=\eps_0(u).
\]
\end{lemma}

\begin{proof}
The integral in \eqref{eq:iid-pipi} is positive semidefinite, and
$\E[(w_b/(W+w_b))^2]\ge\rho_b^2$ by Jensen's inequality. So
$\E[\pi\pi^{\mathsf T}]\succeq\operatorname{diag}(\lambda_b\rho_b^2,\,0)$. By
Lemma~\ref{lem:iid-mecke},
$\E\operatorname{diag}(\pi)=\operatorname{diag}(\lambda_b\rho_b,\eps_0)$. Hence
\[
\Sigma(u)=\operatorname{diag}(\lambda_b\rho_b,\eps_0)-\E\bigl[\pi\pi^{\mathsf T}\bigr]
\preceq\operatorname{diag}\bigl(\lambda_b\rho_b-\lambda_b\rho_b^2,\ \eps_0\bigr).\qedhere
\]
\end{proof}

\begin{lemma}[Legendre facts]\label{lem:iid-legendre}
Let $\mathcal{D}=\nabla\Phi(\R^{B+1})=\nabla\Phi(\mathcal{H}_0)$.
\begin{itemize}
\item[(i)] $\nabla\Phi$ is a $C^\infty$ diffeomorphism from $\mathcal{H}_0$ onto
$\mathcal{D}$, and $\mathcal{D}$ is a relatively open subset of $\mathcal{H}_1$ contained
in $\{\mu:0<\mu_b<\lambda_b\ (b\le B),\ \eps>0\}$. For $\mu=\mu(u)\in\mathcal{D}$,
\[
H(\mu(u))=\Phi(u)-\langle\mu(u),u\rangle .
\]
On $\mathcal{D}$ the function $H$ is $C^2$; at $\mu=\mu(u)$ its derivative in a direction
$x\in\mathcal{H}_0$ is $-\langle u,x\rangle$ and its second derivative in that direction is
$-x^{\mathsf T}(\Sigma(u)|_{\mathcal{H}_0})^{-1}x$.
\item[(ii)] If $\mu_0\in\mathcal{D}$, $\mu\in\mathcal{H}_1$ and $H(\mu)>-\infty$, then
$\mu_t=\mu_0+t(\mu-\mu_0)\in\mathcal{D}$ for every $t\in[0,1)$.
\end{itemize}
\end{lemma}

\begin{proof}
(i) The two descriptions of $\mathcal{D}$ agree because $\nabla\Phi(u+t\ind)=\nabla\Phi(u)$.
The restriction of $\Phi$ to $\mathcal{H}_0$ is smooth and strictly convex, so for
$u\ne u'$ in $\mathcal{H}_0$,
$\langle\nabla\Phi(u)-\nabla\Phi(u'),u-u'\rangle>0$; thus $\nabla\Phi$ is injective on
$\mathcal{H}_0$. Its derivative at $u$, as a map from $\mathcal{H}_0$ to the direction
space $\mathcal{H}_0$ of $\mathcal{H}_1$, is $\Sigma(u)|_{\mathcal{H}_0}$, which is
invertible. By the inverse function theorem $\nabla\Phi$ is a local diffeomorphism from
$\mathcal{H}_0$ to $\mathcal{H}_1$; being injective, it is a diffeomorphism onto its
image, which is open in $\mathcal{H}_1$. The inclusion is Lemma~\ref{lem:iid-mecke}.

Let $\mu=\mu(u)$. The convex function $u'\mapsto\Phi(u')-\langle\mu,u'\rangle$ on
$\R^{B+1}$ has gradient $\nabla\Phi(u)-\mu=0$ at $u$, so it is minimised there, and
$H(\mu)=\Phi(u)-\langle\mu,u\rangle$. Write $u=u(\mu)\in\mathcal{H}_0$ for the inverse
map, which is smooth. Differentiating $H(\mu)=\Phi(u(\mu))-\langle\mu,u(\mu)\rangle$ in a
direction $x\in\mathcal{H}_0$ gives
$\langle\nabla\Phi(u)-\mu,\partial_xu\rangle-\langle x,u\rangle=-\langle u,x\rangle$.
Differentiating once more gives $-\langle\partial_xu,x\rangle
=-x^{\mathsf T}(\Sigma(u)|_{\mathcal{H}_0})^{-1}x$, because the derivative of
$\mu\mapsto u(\mu)$ is the inverse of $\Sigma(u)|_{\mathcal{H}_0}$.

(ii) Let $\varphi_0(u)=\Phi(u)-\langle\mu_0,u\rangle$ for $u\in\mathcal{H}_0$. It is
strictly convex and, by (i), has a minimiser $u_0$, which is therefore unique. Let
$a>0$ be the minimum of $\varphi_0-\varphi_0(u_0)$ on the unit sphere of $\mathcal{H}_0$
centred at $u_0$; by convexity, $\varphi_0(u)-\varphi_0(u_0)\ge a\,|u-u_0|$ whenever
$|u-u_0|\ge1$. Fix $t\in[0,1)$. For $u\in\mathcal{H}_0$,
\[
\Phi(u)-\langle\mu_t,u\rangle=(1-t)\,\varphi_0(u)+t\bigl[\Phi(u)-\langle\mu,u\rangle\bigr]
\ge(1-t)\,\varphi_0(u)+t\,H(\mu),
\]
which tends to $+\infty$ as $|u|\to\infty$ in $\mathcal{H}_0$. So the continuous function
$u\mapsto\Phi(u)-\langle\mu_t,u\rangle$ attains its minimum over $\mathcal{H}_0$ at some
$u_t$. There the vector $\nabla\Phi(u_t)-\mu_t$ is orthogonal to $\mathcal{H}_0$; it also
lies in $\mathcal{H}_0$, because both $\nabla\Phi(u_t)$ and $\mu_t$ lie in $\mathcal{H}_1$.
Hence $\nabla\Phi(u_t)=\mu_t$ and $\mu_t\in\mathcal{D}$.
\end{proof}

\begin{proposition}[Tangent inequality]\label{prop:iid-tangent}
Let $u_0\in\R^{B+1}$ and write $\rho_b=\rho_b(u_0)$, $\eps_0=\eps_0(u_0)$. For every
$\mu=(\mu_1,\dots,\mu_B,\eps)\in\mathcal{P}$,
\begin{multline}\label{eq:iid-tangent}
H(\mu)-S_0(\mu)\le\Phi(u_0)+\sum_{b=1}^B\lambda_b\ln(1-\rho_b)-\eps_0\\
-\sum_{b=1}^B\mu_b\Bigl[u_{0,b}+\ln\frac{1-\rho_b}{\rho_b}\Bigr]
-\eps\bigl[u_{0,\infty}-\ln\eps_0\bigr].
\end{multline}
\end{proposition}

\begin{proof}
Both sides of \eqref{eq:iid-tangent} are unchanged when $u_0$ is replaced by $u_0+t\ind$
(the terms $\Phi(u_0)$ and $-\langle\mu,u_0\rangle$ change by $t$ and $-t$), so one may
assume $u_0\in\mathcal{H}_0$. If $H(\mu)=-\infty$ there is nothing to prove. Otherwise let
$\mu_0=\mu(u_0)\in\mathcal{D}$, $x=\mu-\mu_0\in\mathcal{H}_0$,
$\mu_t=\mu_0+tx$ and $\varphi(t)=H(\mu_t)$ for $t\in[0,1]$. The function $\varphi$ is
concave and finite at both ends, hence finite on $[0,1]$. By
Lemma~\ref{lem:iid-legendre}(ii), $\mu_t\in\mathcal{D}$ for $t\in[0,1)$; let
$u_t\in\mathcal{H}_0$ be the point with $\mu(u_t)=\mu_t$. Since $\mathcal{D}\subset
\{0<\mu_b<\lambda_b,\ \eps>0\}$ and $\mathcal{P}$ is convex, $\mu_t\in\mathcal{P}$ for all
$t\in[0,1]$, and both $\varphi$ and $t\mapsto S_0(\mu_t)$ are $C^2$ on $[0,1)$.

\emph{Second derivatives.} By Lemma~\ref{lem:iid-legendre}(i),
$\varphi''(t)=-x^{\mathsf T}(\Sigma(u_t)|_{\mathcal{H}_0})^{-1}x$. The function $S_0$ is a
sum of functions of the single coordinates, with
\begin{align*}
\frac{\partial S_0}{\partial\mu_b}&=\ln\frac{1-\mu_b/\lambda_b}{\mu_b/\lambda_b},&
\frac{\partial^2S_0}{\partial\mu_b^2}&=-\frac{1}{\lambda_b\,(\mu_b/\lambda_b)(1-\mu_b/\lambda_b)},\\
\frac{\partial S_0}{\partial\eps}&=-\ln\eps,&
\frac{\partial^2S_0}{\partial\eps^2}&=-\frac1\eps .
\end{align*}
At $\mu_t=\mu(u_t)$ one has $\mu_{t,b}/\lambda_b=\rho_b(u_t)$ and
$\eps_t=\eps_0(u_t)$ by Lemma~\ref{lem:iid-mecke}, so the Hessian of $S_0$ at $\mu_t$ is
$-\Gamma(u_t)^{-1}$ and
$\frac{\mathrm{d}^2}{\mathrm{d}t^2}S_0(\mu_t)=-x^{\mathsf T}\Gamma(u_t)^{-1}x$.

\emph{The matrix inequality.} Write $\Sigma=\Sigma(u_t)$, $\Gamma=\Gamma(u_t)$. Since
$x\in\mathcal{H}_0$ and $\Sigma\ind=0$,
\[
x^{\mathsf T}(\Sigma|_{\mathcal{H}_0})^{-1}x
=\sup_{y\in\R^{B+1}}\bigl[2\langle x,y\rangle-y^{\mathsf T}\Sigma y\bigr]
\ge\sup_{y\in\R^{B+1}}\bigl[2\langle x,y\rangle-y^{\mathsf T}\Gamma y\bigr]
=x^{\mathsf T}\Gamma^{-1}x .
\]
The first equality holds because the bracket depends on $y$ only through its projection to
$\mathcal{H}_0$, and on $\mathcal{H}_0$ the supremum of
$2\langle x,y\rangle-y^{\mathsf T}\Sigma y$ is attained at
$y=(\Sigma|_{\mathcal{H}_0})^{-1}x$. The inequality is Lemma~\ref{lem:iid-domination}. The
last equality holds because $\Gamma$ is diagonal with positive entries.

\emph{The tangent at $t=0$.} Consequently
$t\mapsto\varphi(t)-S_0(\mu_t)$ has nonpositive second derivative on $[0,1)$, so it is
concave there and lies below its tangent at $t=0$. By Lemma~\ref{lem:iid-legendre}(i),
$\varphi'(0)=-\langle u_0,x\rangle$. Hence for $t\in[0,1)$
\begin{equation}\label{eq:iid-tangent-t}
\varphi(t)-S_0(\mu_t)\le H(\mu_0)-S_0(\mu_0)-t\,\bigl\langle u_0+\nabla S_0(\mu_0),\,x\bigr\rangle,
\end{equation}
where $\nabla S_0(\mu_0)$ has coordinates $\ln((1-\rho_b)/\rho_b)$ for $b\le B$ and
$-\ln\eps_0$.

\emph{The boundary.} Let $t\uparrow1$ in \eqref{eq:iid-tangent-t}. The function $S_0$ is
continuous on $\mathcal{P}$, so $S_0(\mu_t)\to S_0(\mu)$. Since $\varphi$ is concave and
finite on $[0,1]$, $\varphi(t)\ge(1-t)\varphi(0)+t\varphi(1)$, so
$\liminf_{t\uparrow1}\varphi(t)\ge\varphi(1)=H(\mu)$. Therefore
\[
H(\mu)-S_0(\mu)\le H(\mu_0)-S_0(\mu_0)-\bigl\langle u_0+\nabla S_0(\mu_0),\,\mu-\mu_0\bigr\rangle .
\]

\emph{Evaluation.} By Lemma~\ref{lem:iid-legendre}(i),
$H(\mu_0)+\langle u_0,\mu_0\rangle=\Phi(u_0)$. Using $h(r)-r\,h'(r)=-\ln(1-r)$ and
$\mu_{0,b}=\lambda_b\rho_b$,
\begin{multline*}
S_0(\mu_0)-\langle\nabla S_0(\mu_0),\mu_0\rangle
=\sum_b\lambda_b\bigl[h(\rho_b)-\rho_bh'(\rho_b)\bigr]
 +\bigl[-\eps_0\ln\eps_0+\eps_0+\eps_0\ln\eps_0\bigr]\\
=-\sum_b\lambda_b\ln(1-\rho_b)+\eps_0 .
\end{multline*}
Substituting these two identities gives \eqref{eq:iid-tangent}.
\end{proof}

\begin{proposition}[Tangent bound]\label{prop:iid-tangent-bound}
For every $K>0$, $B\ge1$, $\theta>0$, $\beta\in\R$ and $u_0\in\R^{B+1}$,
\[
L^\ast=\sup_{\mu\in\mathcal{P}}L(\mu)\le\theta\beta+\Psi_0(K,B,\theta,u_0).
\]
\end{proposition}

\begin{proof}
Write $L(\mu)-\theta\beta=3\,(H(\mu)-S_0(\mu))+S_0(\mu)+\langle\ell,\mu\rangle$ and insert
\eqref{eq:iid-tangent}:
\begin{equation}\label{eq:iid-after-tangent}
L(\mu)-\theta\beta\le3\Bigl[\Phi(u_0)+\sum_b\lambda_b\ln(1-\rho_b)-\eps_0\Bigr]
+S_0(\mu)+\langle a,\mu\rangle,
\end{equation}
with
\[
a_b=\ell_b-3u_{0,b}-3\ln\frac{1-\rho_b}{\rho_b}\quad(b\le B),\qquad
a_\infty=\ell_\infty-3u_{0,\infty}+3\ln\eps_0 .
\]
Now drop the constraint $\sum_b\mu_b+\eps=1$ and maximise coordinate by coordinate. For
real $a$, $\sup_{0\le r\le1}[h(r)+ar]=\ln(1+\mathrm{e}^a)$ and
$\sup_{\eps\ge0}[-\eps\ln\eps+\eps+a\eps]=\mathrm{e}^a$. With $r=\mu_b/\lambda_b$ this
gives
\begin{align*}
S_0(\mu)+\langle a,\mu\rangle
&=\sum_b\lambda_b\Bigl[h\Bigl(\frac{\mu_b}{\lambda_b}\Bigr)+a_b\frac{\mu_b}{\lambda_b}\Bigr]
+\bigl[-\eps\ln\eps+\eps+a_\infty\eps\bigr]\\
&\le\sum_b\lambda_b\ln\bigl(1+\mathrm{e}^{a_b}\bigr)+\mathrm{e}^{a_\infty}.
\end{align*}
Finally $\mathrm{e}^{a_b}=C_b(u_0)\,\rho_b^3/(1-\rho_b)^3$, so
$(1-\rho_b)^3(1+\mathrm{e}^{a_b})=(1-\rho_b)^3+C_b\rho_b^3$, and
$\mathrm{e}^{a_\infty}=C_\infty(u_0)\,\eps_0^3$. The right side of
\eqref{eq:iid-after-tangent} is therefore at most
$3\Phi(u_0)+\sum_b\lambda_b\ln((1-\rho_b)^3+C_b\rho_b^3)-3\eps_0+C_\infty\eps_0^3
=\Psi_0(K,B,\theta,u_0)$.
\end{proof}

\begin{remark}\label{rem:iid-binning}
The tilt uses the exact conditional Laplace transform $\kappa_b/\lambda_b$ of the cost
within a bin, see \eqref{eq:iid-laplace}, so the binning causes no error in the
accounting of costs; the only loss is that the conditioning is on coarser stars. Every
$B\ge1$ gives a valid bound. Likewise Proposition~\ref{prop:iid-tangent-bound} holds for
every weight vector $u_0$, so the certificate does not need an exact stationary point: the
vector $u^\ast$ of Proposition~\ref{prop:iid-certificate} was found by a floating-point
search, and how it was found is irrelevant to the proof.
\end{remark}

\subsection{Proof of the criterion}\label{subsec:iid-proof}

\begin{proof}[Proof of Proposition~\ref{prop:iid-criterion}]
Let $K,B,\theta,u,\beta$ satisfy \eqref{eq:iid-criterion}. By
Proposition~\ref{prop:iid-tangent-bound} with $u_0=u$,
$L^\ast\le\theta\beta+\Psi_0(K,B,\theta,u)<0$. By Lemma~\ref{lem:iid-limit} there are
$\delta>0$ and $c>0$, depending only on $(K,B,\theta,u,\beta)$, such that
$\E[Z_\beta;D(G)\in\mathcal{T}_\delta]\le\mathrm{e}^{-cq}$ for all large $q$. By Markov's
inequality,
\[
\P(Z_\beta\ge1)\le\P\bigl(D(G)\notin\mathcal{T}_\delta\bigr)
+\E\bigl[Z_\beta;\,D(G)\in\mathcal{T}_\delta\bigr]
\le\P\bigl(D(G)\notin\mathcal{T}_\delta\bigr)+\mathrm{e}^{-cq},
\]
and the first term tends to $0$ by Lemma~\ref{lem:iid-typical}.
\end{proof}

\begin{remark}[The certificate]\label{rem:iid-certificate}
For the parameters of Proposition~\ref{prop:iid-certificate} the script reports the
following outward-rounded enclosures, which are intermediate quantities and are not used
elsewhere:
\begin{align*}
\Phi(u^\ast)&\in[-0.522226313363,\,-0.522223198517],\\
\eps_0(u^\ast)&\in[0.02081077388,\,0.02081090331],\\
\rho_1(u^\ast)&\in[0.7564957391,\,0.7565004438],\\
\textstyle\sum_b\lambda_b\rho_b(u^\ast)+\eps_0(u^\ast)&\in[0.99999792704,\,1.00000414596].
\end{align*}
The last interval contains $1$, as it must by Lemma~\ref{lem:iid-mecke}. The certified
upper bounds for $\theta\beta+\Psi_0$ are $-0.0037532832$ at $\beta=2.30$,
$-0.0004108832$ at $\beta=2.302$, $-0.0000766432$ at $\beta=2.3022$ and $-0.0000097952$
at $\beta=2.30224$. The enclosure of $\Psi_0$ is rigorous given the correctness of the
script, of exact integer and rational arithmetic, and of the four interval operations and
the comparisons of the interval library (mpmath~1.3.0). The library's interval exponential
and logarithm are not used, because they are not strict enclosures; see
Section~\ref{sec:certificates}.
\end{remark}

\begin{remark}[Checks that carry no weight in the proof]\label{rem:iid-checks}
Lemma~\ref{lem:iid-count} was compared with an exact enumeration of all $20{,}736$
configurations of an instance with three agents per sector and two bins, for all ten
profiles (\texttt{lemmaC/\allowbreak round3/\allowbreak firstmoment/\allowbreak
check\_partial.py}).
Lemma~\ref{lem:iid-domination} and inequality \eqref{eq:iid-tangent} were tested in
floating point on random instances with one to three bins
(\texttt{lemmaC/\allowbreak round3/\allowbreak firstmoment/\allowbreak
check\_concave.py}): the largest eigenvalue of
$\Gamma^{-1/2}\Sigma\,\Gamma^{-1/2}$ found over $60$ instances was $0.999999$, and the
largest value of the left side of \eqref{eq:iid-tangent} minus its right side over $360$
pairs $(u_0,\mu)$ was $-1.3\times10^{-11}$.
\end{remark}
